Este trabalho aborda o reconhecimento de dígitos manuscritos no MNIST (70.000 imagens $28 \times 28$, 10 classes) por meio de um MLP de duas camadas ($784 \to h \to 10$) em NumPy (\texttt{learn.py}), com ativação ReLU, saída Softmax numericamente estável, perda de entropia cruzada e treinamento via gradiente descendente em lote completo (\emph{full-batch GD}). A partir da base ($\alpha = 0{,}1$, 200 épocas, $h = 64$), variaram-se controladamente a taxa de aprendizado ($\alpha \in \{0{,}001;\; 0{,}01;\; 0{,}1;\; 0{,}5\}$), as épocas ($\{50;\; 100;\; 200;\; 500\}$) e os neurônios escondidos ($h \in \{4;\; 16;\; 64;\; 256\}$). O melhor modelo ($\alpha = 0{,}5$, 200 épocas, $h = 64$) atingiu 92,43\% de acurácia de teste (perda de treino $0{,}2585$), superando o limiar de 90\%. Os resultados evidenciam forte sensibilidade à taxa de aprendizado e compromissos entre capacidade, tempo de treino e generalização.
